\chapter{Mod-\texorpdfstring{$\varphi$}{phi} Convergence}
\begin{center}\itshape
The central limit theorem is a magnifying glass that shows only the bell curve. Divide the characteristic function by the bell curve\\
before taking the limit, and something survives: a ``residue'' that remembers where the sum came from.\\
The cycles of a permutation and the prime factors of an integer leave the same fingerprint, $1/\Gamma(z)$.
\end{center}

By the Erd\H{o}s--Kac theorem of Volume~5, the number $\omega(n)$ of distinct prime factors of a random integer $n\le x$ is approximately Gaussian, with mean and variance $\log\log x$. The theorem is true,
but almost useless in practice: $\log\log x$ grows so slowly that at $x=10^{100}$ it is only $5.4$, and the error term is of order $1/\sqrt{\log\log x}$, as R\'enyi and Tur\'an showed. The Gaussian limit has thrown
away almost everything. What does it throw away, and can we keep it?

The answer came from an unexpected direction. In 2000, Jon Keating and Nina Snaith computed the characteristic function of $\log|\det(I-U)|$ for a random unitary matrix $U$ of size $N$, a model for $\log|\zeta(\frac12+it)|$.
It is Gaussian with variance $\frac12\log N$, up to a factor which converges, as $N\to\infty$, to an explicit product of Barnes functions. Their conjecture for the moments of the zeta function, built on that factor, matched every
known result of number theory. A decade later Jean Jacod, Emmanuel Kowalski and Ashkan Nikeghbali made the factor into a definition: a sequence \emph{converges mod-Gaussian} if its characteristic function, divided by the
Gaussian one, has a limit. Kowalski and Nikeghbali did the same with Poisson laws, and the theory of \emph{mod-$\varphi$ convergence} was born, with Valentin F\'eray and Pierre-Lo\"ic M\'eliot developing its consequences.

This chapter proves the two basic consequences of mod-Poisson convergence, a central limit theorem and a local limit theorem at the scale of large deviations, in which the limiting factor is the correction to the
Poisson law. Then we compute the factor for the cycles of a random permutation, where it is $1/\Gamma(z)$, and for the prime factors of a random integer, where it is an arithmetic Euler product, and we see how the two fit together.

\begin{reminder}[from Volumes 3, 4 and 5, and Chapter 3]
The number of cycles of a uniform random permutation of $\{1,\dots,n\}$ has the law of $B_1+\dots+B_n$, independent with $B_j$ Bernoulli$(\frac1j)$ (Feller's coupling, Volume~3, Theorem~2.12). L\'evy's continuity theorem (Theorem~\ref{thm:levycont}):
convergence of characteristic functions to a characteristic function implies convergence in law. Weierstrass' product $\frac1{\Gamma(x)}=xe^{\gamma x}\prod_{k\ge1}(1+\frac xk)e^{-x/k}$ for $x>0$ (Volume~4, Corollary~5.7); $\frac1\Gamma$ extends to an entire function
(Volume~4, Theorem~9.36), and two holomorphic functions on a connected open set that agree on a segment are equal (Volume~4, Theorem~9.18). $\sum_p\frac1p$ diverges, and $\sum_{p\le x}\frac1p=\log\log x+M+o(1)$ (Mertens, Volume~5, Theorem~1.5).
\end{reminder}

\section{Beyond the central limit theorem}

\begin{definition}\label{def:modgauss}
Real random variables $X_n$ \emph{converge mod-Gaussian}, with parameters $\sigma_n^2\to\infty$ and limiting function $\psi$, if for every real $t$
\[
\mathbb E\bigl[e^{itX_n}\bigr]\,e^{\sigma_n^2t^2/2}\longrightarrow\psi(t),
\]
locally uniformly in $t$, where $\psi$ is continuous with $\psi(0)=1$.
\end{definition}

If $\psi$ were the characteristic function of a variable $Y$, the definition would say that $X_n$ looks like a Gaussian of variance $\sigma_n^2$ \emph{plus} an independent copy of $Y$. In general $\psi$ need not be a characteristic function, and
it is better thought of as a residue: what is left of the law of $X_n$ after the Gaussian part has been divided out.

\begin{proposition}\label{prop:modgaussclt}
If $X_n$ converges mod-Gaussian with parameters $\sigma_n^2$, then $X_n/\sigma_n\to\mathcal N(0,1)$ in law.
\end{proposition}
\begin{proof}
For fixed $t$, $\mathbb Ee^{itX_n/\sigma_n}=e^{-t^2/2}\psi_n(t/\sigma_n)$, where $\psi_n(u)=\mathbb E[e^{iuX_n}]e^{\sigma_n^2u^2/2}$ tends to $\psi$ uniformly on $[-1,1]$. Since $t/\sigma_n\to0$ and $\psi$ is continuous with $\psi(0)=1$, $\psi_n(t/\sigma_n)\to1$. Apply L\'evy's continuity theorem.
\end{proof}

\begin{example}[a harmonic random walk]\label{ex:harmonicwalk}
Let $\varepsilon_k$ be independent fair signs and $X_n=\sum_{k\le n}\varepsilon_k/\sqrt k$, with variance $\sigma_n^2=H_n=\sum_{k\le n}\frac1k$. Then
\[
\mathbb E\bigl[e^{itX_n}\bigr]e^{H_nt^2/2}=\prod_{k\le n}\cos\Bigl(\frac t{\sqrt k}\Bigr)e^{t^2/2k}.
\]
For $|u|\le1$, $\log\cos u=-\frac{u^2}2-\frac{u^4}{12}+O(u^6)$, so the $k$-th factor is $\exp(-\frac{t^4}{12k^2}+O(\frac{t^6}{k^3}))$ once $k\ge t^2$, and the product converges locally uniformly to $\psi(t)=\prod_{k\ge1}\cos(t/\sqrt k)e^{t^2/2k}$:
mod-Gaussian convergence. Replacing the signs by other symmetric variables of variance $1$ changes $\psi$ but not the Gaussian limit (Exercise~\ref{ex:residues}). The residue remembers what the central limit theorem forgets.
\end{example}

For integer-valued variables, such as counts, the natural reference law is the Poisson law, whose generating function $e^{\lambda(z-1)}$ makes sense for all complex $z$.

\begin{definition}\label{def:modpoisson}
Let $X_n$ be random variables with values in $\{0,1,2,\dots\}$, $\lambda_n\to\infty$, and $0<r_1<1<r_2\le\infty$. The $X_n$ \emph{converge mod-Poisson} on the annulus $A=\{r_1<|z|<r_2\}$, with parameters $\lambda_n$ and limiting function
$\psi$, if $\mathbb E[r_2'^{\,X_n}]<\infty$ for every $r_2'<r_2$, and
\[
\psi_n(z)=e^{-\lambda_n(z-1)}\,\mathbb E\bigl[z^{X_n}\bigr]\longrightarrow\psi(z)
\]
uniformly on every closed annulus $\{a\le|z|\le b\}$ with $r_1<a\le b<r_2$, where $\psi$ is continuous on $A$ and $\psi(1)=1$.
\end{definition}

\begin{example}[sums of rare independent events]\label{ex:bernsums}
Let $X_n=\sum_{j\le n}B_j$ with independent Bernoulli$(p_j)$ variables, $\lambda_n=\sum_{j\le n}p_j\to\infty$ and $\sum_jp_j^2<\infty$. Then
\[
\psi_n(z)=\prod_{j\le n}\bigl(1+(z-1)p_j\bigr)e^{-(z-1)p_j}.
\]
With $w=z-1$ and $|wp_j|\le\frac12$, the factor is $\exp(\log(1+wp_j)-wp_j)$ with $|\log(1+u)-u|\le|u|^2$ for $|u|\le\frac12$. So on $|z|\le b$ the tail of the product beyond the finitely many $j$ with $(b+1)p_j>\frac12$ converges
uniformly, and $\psi_n\to\psi(z)=\prod_{j\ge1}(1+(z-1)p_j)e^{-(z-1)p_j}$ uniformly on compact subsets of $\C$: mod-Poisson convergence on $A=\C\setminus\{0\}$ ($r_1=0$, $r_2=\infty$; we may take $r_1$ as small as we like). The limit $\psi$ is an entire function vanishing
exactly at the points $1-\frac1{p_j}$.
\end{example}

The Poisson law itself converges mod-Poisson with $\psi=1$, and adding an independent bounded variable $Y$ multiplies $\psi$ by $\mathbb E[z^Y]$ (Exercise~\ref{ex:shift}).

\section{What mod-Poisson convergence gives}

\begin{theorem}[central limit theorem]\label{thm:modpoissonclt}
If $X_n$ converges mod-Poisson with parameters $\lambda_n$, then $\frac{X_n-\lambda_n}{\sqrt{\lambda_n}}\to\mathcal N(0,1)$ in law.
\end{theorem}
\begin{proof}
For real $t$, put $z_n=e^{it/\sqrt{\lambda_n}}$, on the unit circle, inside $A$. Then
\[
\mathbb E\bigl[e^{it(X_n-\lambda_n)/\sqrt{\lambda_n}}\bigr]=e^{-it\sqrt{\lambda_n}}e^{\lambda_n(z_n-1)}\psi_n(z_n)=\exp\Bigl(\lambda_n\bigl(e^{it/\sqrt{\lambda_n}}-1-\tfrac{it}{\sqrt{\lambda_n}}\bigr)\Bigr)\psi_n(z_n).
\]
The exponent is $\lambda_n(-\frac{t^2}{2\lambda_n}+O(\lambda_n^{-3/2}))\to-\frac{t^2}2$. And $\psi_n(z_n)-\psi(z_n)\to0$ by uniform convergence on $|z|=1$, while $\psi(z_n)\to\psi(1)=1$ by continuity. L\'evy's continuity theorem concludes.
\end{proof}

The real content of mod-Poisson convergence is at the scale of large deviations, where the central limit theorem says nothing. Write $\pi_\lambda(k)=e^{-\lambda}\frac{\lambda^k}{k!}$ for the Poisson probabilities.

\begin{theorem}[local limit theorem with residue]\label{thm:modpoissonlocal}
Let $X_n$ converge mod-Poisson on $A=\{r_1<|z|<r_2\}$ with parameters $\lambda_n$ and limit $\psi$, and let $r_1<a\le b<r_2$. Then, uniformly in the integers $k$ with $a\le\frac k{\lambda_n}\le b$,
\[
P(X_n=k)=\pi_{\lambda_n}(k)\Bigl(\psi\Bigl(\frac k{\lambda_n}\Bigr)+o(1)\Bigr)\qquad(n\to\infty).
\]
\end{theorem}
\begin{proof}
Write $\lambda=\lambda_n$, $\varphi_n(z)=\mathbb E[z^{X_n}]$, and let $r=k/\lambda\in[a,b]$. Since $\sum_jP(X_n=j)r^j<\infty$, we may integrate the series term by term, and orthogonality of the $e^{ij\theta}$ gives
\[
\begin{gathered}P(X_n=k)=\frac1{2\pi}\int_{-\pi}^\pi\varphi_n(re^{i\theta})\,r^{-k}e^{-ik\theta}\,\dd\theta,\qq\pi_\lambda(k)=\frac1{2\pi}\int_{-\pi}^\pi e^{\lambda(re^{i\theta}-1)}r^{-k}e^{-ik\theta}\,\dd\theta .\end{gathered}
\]
Since $\varphi_n(z)=e^{\lambda(z-1)}\psi_n(z)$ and
\[
\bigl|e^{\lambda(re^{i\theta}-1)}\bigr|=e^{\lambda(r-1)}e^{-\lambda r(1-\cos\theta)},
\]
we get
\[
\begin{aligned}\bigl|P(X_n=k)-\psi(r)\pi_\lambda(k)\bigr|&\le\frac{e^{\lambda(r-1)}r^{-k}}{2\pi}\eqbr\times\int_{-\pi}^\pi e^{-\lambda r(1-\cos\theta)}\bigl|\psi_n(re^{i\theta})-\psi(r)\bigr|\,\dd\theta .\end{aligned}
\]
Let $\varepsilon>0$ and $M=1+\sup_{a\le|z|\le b}|\psi(z)|$. For $n$ large, $|\psi_n-\psi|\le\varepsilon$ on the closed annulus, and $|\psi_n|\le M$. By uniform continuity of $\psi$ there, there is $\delta>0$ with $|\psi(re^{i\theta})-\psi(r)|\le\varepsilon$ for $|\theta|\le\delta$ and all
$r\in[a,b]$. On $[-\pi,\pi]$, $1-\cos\theta\ge\frac{2\theta^2}{\pi^2}$. Hence, using $\lambda r=k$, the integral is at most
\[
2\varepsilon\int_\R e^{-2k\theta^2/\pi^2}\,\dd\theta+4\pi Me^{-2a\lambda\delta^2/\pi^2}=2\varepsilon\sqrt{\frac{\pi^3}{2k}}+4\pi Me^{-2a\lambda\delta^2/\pi^2}.
\]
On the other hand, $k!\le e\sqrt k\,(k/e)^k$ for $k\ge1$ (Volume~4, Theorem~5.8 gives more), so
\[
\pi_\lambda(k)\ge\frac{e^{-\lambda}\lambda^ke^k}{e\sqrt k\,k^k}=\frac{e^{\lambda(r-1)}r^{-k}}{e\sqrt k}.
\]
Dividing,
\[
\Bigl|\frac{P(X_n=k)}{\pi_\lambda(k)}-\psi(r)\Bigr|\le e\,\varepsilon\sqrt{\frac\pi2}+2eM\sqrt{b\lambda}\,e^{-2a\lambda\delta^2/\pi^2},
\]
and the last term tends to $0$ as $n\to\infty$, uniformly in $k$. Since $\varepsilon$ is arbitrary, the theorem follows.
\end{proof}

At $k\approx\lambda_n$, where $\psi(1)=1$, the theorem says that $X_n$ is locally Poisson, a strong form of the central limit theorem. At $k=x\lambda_n$ with $x\neq1$, it says that the probability of a large deviation is the Poisson one
multiplied by $\psi(x)$: the residue, invisible to the central limit theorem, is exactly the correction factor for rare events. The bound $k!\le e\sqrt k(k/e)^k$ follows by comparing $\log k!=\sum_{j\le k}\log j$ with
$\int_1^k\log t\,\dd t+\frac12\log k$, the trapezoid rule for a concave function.

\section{Cycles of a random permutation}

Let $K_n$ be the number of cycles of a uniform random permutation of $\{1,\dots,n\}$. Vladimir Goncharov proved in 1944 that $K_n$ is asymptotically Gaussian with mean and variance $\log n$ (Volume~3, Exercise~6.7, asked you to prove it).
Mod-Poisson convergence tells the whole story.

\begin{theorem}\label{thm:cyclesmod}
For every complex $z$,
\[
\mathbb E\bigl[z^{K_n}\bigr]=\prod_{j=1}^n\Bigl(1+\frac{z-1}j\Bigr)=\frac{z(z+1)\cdots(z+n-1)}{n!},
\]
and $K_n$ converges mod-Poisson on $\C\setminus\{0\}$ with parameters $H_n=\sum_{j\le n}\frac1j$ and limiting function
\[
\psi(z)=\prod_{j\ge1}\Bigl(1+\frac{z-1}j\Bigr)e^{-(z-1)/j}=\frac{e^{-\gamma(z-1)}}{\Gamma(z)}.
\]
Equivalently, with parameters $\log n$, the limiting function is $\frac1{\Gamma(z)}$.
\end{theorem}
\begin{proof}
By Feller's coupling, $K_n$ is a sum of independent Bernoulli$(\frac1j)$, $j\le n$, so $\mathbb Ez^{K_n}=\prod_j(1-\frac1j+\frac zj)$, which is the product formula. Example~\ref{ex:bernsums} with $p_j=\frac1j$ gives the mod-Poisson
convergence with limit $\psi(z)=\prod_j(1+\frac wj)e^{-w/j}$, $w=z-1$, an entire function. For real $w>0$, Weierstrass' product at $x=w$ gives
\[
\prod_{j\ge1}\Bigl(1+\frac wj\Bigr)e^{-w/j}=\frac{e^{-\gamma w}}{w\Gamma(w)}=\frac{e^{-\gamma w}}{\Gamma(w+1)}=\frac{e^{-\gamma(z-1)}}{\Gamma(z)}.
\]
Both sides are entire in $z$ and agree for real $z>1$, so they agree everywhere (Volume~4, Theorem~9.18). Finally $H_n-\log n\to\gamma$, so $e^{-(z-1)\log n}\mathbb Ez^{K_n}=e^{(z-1)(H_n-\log n)}\psi_n(z)\to e^{\gamma(z-1)}\psi(z)=\frac1{\Gamma(z)}$, uniformly on compact sets.
\end{proof}

\begin{corollary}\label{cor:cyclesdev}
Uniformly for $a\le\frac k{\log n}\le b$, with $0<a<b$ fixed,
\[
P(K_n=k)=e^{-H_n}\frac{H_n^k}{k!}\Bigl(\frac{e^{-\gamma(x-1)}}{\Gamma(x)}+o(1)\Bigr),\qquad x=\frac k{H_n}.
\]
\end{corollary}
\begin{proof}
Theorem~\ref{thm:modpoissonlocal} with $\lambda_n=H_n$; the ranges $a\le\frac k{\log n}\le b$ and $a'\le\frac k{H_n}\le b'$ differ by a factor tending to $1$.
\end{proof}

So a permutation of a million elements, which has on average $H_{10^6}\approx14.4$ cycles, has twice that many with a probability about $\frac{e^{-\gamma}}{\Gamma(2)}\approx0.56$ times the Poisson prediction, and half that many with
about $\frac{e^{\gamma/2}}{\Gamma(1/2)}\approx0.75$ times it. Figure~\ref{fig:cyclesratio} shows the ratio for $n=10^3$ and $n=10^6$. Since $P(K_n=k)=\frac{|s(n,k)|}{n!}$, where $s(n,k)$ are the Stirling numbers of the first kind, the corollary is
also an asymptotic formula for them (Exercise~\ref{ex:stirling1}).

\begin{figure}[ht]
\centering\fig{ch11-cycles}
\caption{The ratio $P(K_n=k)/\pi_{H_n}(k)$ as a function of $x=k/H_n$, computed exactly, for $n=10^3$ (blue dots) and $n=10^6$ (green dots), and the limiting residue $e^{-\gamma(x-1)}/\Gamma(x)$ (red). The residue equals $1$ at $x=1$: the
Poisson law is right in the bulk, and wrong by the factor $\psi(x)$ in the tails. At finite $n$ the ratio still differs from the limit by a few per cent, an error of order $1/\log n$.}
\label{fig:cyclesratio}
\end{figure}

\section{Prime factors of a random integer}

Choose an integer $N$ at random with $P(N=n)=\frac{n^{-s}}{\zeta(s)}$, for a parameter $s>1$. This ``Dirichlet random integer'' has a marvellous property.

\begin{proposition}\label{prop:dirichletint}
For the Dirichlet random integer $N_s$, the events $\{p\mid N_s\}$, $p$ prime, are independent, with $P(p\mid N_s)=p^{-s}$. Consequently $\omega(N_s)=\sum_pB_p$ with independent Bernoulli$(p^{-s})$ variables $B_p$.
\end{proposition}
\begin{proof}
For $d\ge1$, $P(d\mid N_s)=\sum_{m\ge1}\frac{(dm)^{-s}}{\zeta(s)}=d^{-s}$. For distinct primes $p_1,\dots,p_r$, all divide $N_s$ exactly when their product does, so $P(p_1\mid N_s,\dots,p_r\mid N_s)=(p_1\cdots p_r)^{-s}=\prod_iP(p_i\mid N_s)$. This is
the product rule for every finite subfamily, hence independence of the events (Volume~3: the product rule for intersections implies it for complements). $\omega(N_s)$ counts the primes dividing $N_s$, and it is finite since $N_s$ is.
\end{proof}

The exact independence of divisibility by different primes, which is only approximate for a uniform $n\le x$ and which Volume~5 had to fight for in the Erd\H{o}s--Kac theorem, is free in this model. Euler's product
$\zeta(s)=\prod_p(1-p^{-s})^{-1}$ is the same fact in disguise.

\begin{theorem}\label{thm:omegamod}
Let $\lambda_s=\sum_pp^{-s}$. As $s\downarrow1$, $\lambda_s\to\infty$, and $\omega(N_s)$ converges mod-Poisson on $\C\setminus\{0\}$ with parameters $\lambda_s$ and limiting function
\[
\Psi(z)=\prod_p\Bigl(1+\frac{z-1}p\Bigr)e^{-(z-1)/p}.
\]
\end{theorem}
\begin{proof}
By monotone convergence $\lambda_s\to\sum_p\frac1p=\infty$. By Proposition~\ref{prop:dirichletint},
\[
e^{-\lambda_s(z-1)}\mathbb Ez^{\omega(N_s)}=\prod_p(1+wp^{-s})e^{-wp^{-s}}
\]
with $w=z-1$. Fix $R$ and let $|w|\le R$. For $p>2R$ and every $s\ge1$,
$|wp^{-s}|\le\frac12$ and the logarithm of the $p$-th factor is at most $|w|^2p^{-2}$ in modulus, as in Example~\ref{ex:bernsums}: the tails of the products are uniformly small, uniformly in $s\in[1,2]$. Each of the finitely many factors with
$p\le2R$ is continuous in $(s,w)$. So the product converges, as $s\downarrow1$, uniformly for $|w|\le R$, to its value at $s=1$, which is $\Psi(z)$. (Here the parameter is the continuous $s\downarrow1$ instead of $n\to\infty$; the proofs of the
two theorems above apply verbatim.)
\end{proof}

For a uniform integer $n\le x$ the answer is the same, with one extra factor. Atle Selberg in 1954, and Hubert Delange in general in 1959 and 1971, proved by contour integration around the singularity of $\zeta(s)^z$ at $s=1$ that
\begin{equation}\label{eq:selbergdelange}
\begin{aligned}\frac1x\sum_{n\le x}z^{\omega(n)}&=(\log x)^{z-1}\Bigl(\frac{F(z)}{\Gamma(z)}+O\Bigl(\frac1{\log x}\Bigr)\Bigr),\qq F(z)&=\prod_p\Bigl(1+\frac{z-1}p\Bigr)\Bigl(1-\frac1p\Bigr)^{z-1},\end{aligned}
\end{equation}
uniformly for $|z|\le R$. This is the \emph{Selberg--Delange theorem}; its proof belongs to analytic number theory and uses the contour methods of Volume~5. Volume~14 takes up the distribution of $\omega(n)$ that the identity encodes. In the language of this chapter: $\omega(n)$, for $n$ uniform in $[1,x]$,
converges mod-Poisson with parameters $\log\log x$ and limiting function
\[
\Psi_{\mathrm{unif}}(z)=\frac1{\Gamma(z)}\prod_p\Bigl(1+\frac{z-1}p\Bigr)\Bigl(1-\frac1p\Bigr)^{z-1}=\frac{e^{c(z-1)}}{\Gamma(z)}\,\Psi(z),\qquad c=\sum_p\Bigl(\log\Bigl(1-\frac1p\Bigr)+\frac1p\Bigr).
\]
(The last equality is the convergent product $\prod_p(1-\frac1p)^{z-1}e^{(z-1)/p}=e^{c(z-1)}$; the factor $e^{c(z-1)}$ only shifts the parameter $\lambda$ by a constant.) The residue of a uniform random integer is the product of two factors:
an arithmetic one, $\Psi$, which the Dirichlet model already sees and which comes from the independence of primes; and $\frac1{\Gamma(z)}$, \emph{the same as for permutations}, which comes from the sharp cut-off $n\le x$. The Dirichlet weights cut off
smoothly, and their residue has no Gamma factor.

\begin{figure}[ht]
\centering\fig{ch11-omega}
\caption{The ratio of the proportion of $n\le10^8$ with $\omega(n)=k$ to the Poisson probability $\pi_\lambda(k)$, $\lambda=\log\log10^8$, against $x=k/\lambda$ (dots), and the Selberg--Delange residue $\Psi_{\mathrm{unif}}(x)$ (red). The residue gets the
direction right, fewer integers than Poisson in both tails, but at $\lambda\approx2.9$ the local theorem is far from its limit and the true ratio is much more extreme. No $n\le10^8$ has more than $8$ prime factors, since $2\cdot3\cdots23>10^8$.}
\label{fig:omegaratio}
\end{figure}

The parallel between the cycles of a permutation of size $m$ and the prime factors of an integer of size $x$, with $m\leftrightarrow\log x$, goes much further: the sizes of the prime factors, on a logarithmic scale, have the same limit
law as the sizes of the cycles. That is the subject of Chapter~12.

\begin{think}
Theorem~\ref{thm:modpoissonlocal} needs convergence on an annulus. For $\Omega(N_s)$, the number of prime factors counted with multiplicity, $\mathbb Ez^{\Omega(N_s)}=\prod_p\frac{1-p^{-s}}{1-zp^{-s}}$ has a pole at $z=2^s$ (Exercise~\ref{ex:bigomega}). What does that
say about integers with many more than $2\log\log x$ prime factors? (Think of $n=2^km$.)
\end{think}

\section{Exercises}

\begin{exercise}[$\star$]\label{ex:shift}
Show that Poisson variables of parameters $\lambda_n\to\infty$ converge mod-Poisson with $\psi=1$, and that if $Y$ is a bounded variable with values in $\{0,1,2,\dots\}$, independent of $X_n$, then $X_n+Y$ converges mod-Poisson with $\psi(z)=\mathbb E[z^Y]$.
The central limit theorem cannot tell $X_n$ from $X_n+Y$; Theorem~\ref{thm:modpoissonlocal} can.
\end{exercise}

\begin{exercise}[$\star$]\label{ex:lambdas}
Show that for the Dirichlet random integer, $P(N_s\le y)\to0$ as $s\downarrow1$ for every fixed $y$: the model describes ``large'' integers. Show that $\mathbb E\,\omega(N_s)=\lambda_s$ and that $\lambda_s=\log\frac1{s-1}+O(1)$. (Compare $\sum_pp^{-s}$ with
$\log\zeta(s)=\sum_p\sum_m\frac{p^{-ms}}m$ and use $\zeta(s)=\frac1{s-1}+O(1)$.)
\end{exercise}

\begin{exercise}[$\star\star$]\label{ex:residues}
Let $U_k$ be independent and uniform on $[-\sqrt3,\sqrt3]$ (variance $1$) and $Y_n=\sum_{k\le n}U_k/\sqrt k$. Show that $Y_n$ converges mod-Gaussian with parameters $H_n$ and compute its limiting function. Show that it differs from the one of
Example~\ref{ex:harmonicwalk}, by comparing the coefficients of $t^4$ in their logarithms.
\end{exercise}

\begin{exercise}[$\star\star$]\label{ex:bigomega}
Show that $\mathbb Ez^{\Omega(N_s)}=\prod_p\frac{1-p^{-s}}{1-zp^{-s}}$ for $|z|<2^s$, and that $\Omega(N_s)$ converges mod-Poisson on $\{0<|z|<2\}$ with parameters $\lambda_s$ and limiting function
$\prod_p\frac{1-1/p}{1-z/p}e^{-(z-1)/p}$. Why can the annulus not be larger?
\end{exercise}

\begin{exercise}[$\star\star$]\label{ex:stirling1}
The unsigned Stirling numbers of the first kind are defined by $z(z+1)\cdots(z+n-1)=\sum_k|s(n,k)|z^k$. Show that, uniformly for $a\le x\le b$ with $x=\frac k{\log n}$ and $0<a<b$ fixed,
\[
|s(n,k)|\sim\frac{(n-1)!\,(\log n)^{k-1}}{(k-1)!\;\Gamma(1+x)}.
\]
(Use Theorem~\ref{thm:cyclesmod} with parameters $\log n$.)
\end{exercise}

\section{Hints}

\begin{hint}{ex:shift}
$\mathbb Ez^{X_n+Y}=e^{\lambda_n(z-1)}\mathbb Ez^Y$.
\end{hint}
\begin{hint}{ex:lambdas}
$P(N_s\le y)\le\frac{\sum_{n\le y}n^{-1}}{\zeta(s)}$ and $\zeta(s)\ge\int_1^\infty t^{-s}\dd t=\frac1{s-1}$. The terms $m\ge2$ in $\log\zeta(s)$ sum to at most $\sum_pp^{-2}(1-p^{-1})^{-1}<\infty$.
\end{hint}
\begin{hint}{ex:residues}
$\mathbb Ee^{iuU}=\frac{\sin\sqrt3u}{\sqrt3u}$, and $\log\frac{\sin v}v=-\frac{v^2}6-\frac{v^4}{180}+O(v^6)$.
\end{hint}
\begin{hint}{ex:bigomega}
The exponent of $p$ in $N_s$ is geometric: $P(p^m\mid N_s)=p^{-ms}$. At $z=2$, the factor for $p=2$ blows up as $s\downarrow1$.
\end{hint}
\begin{hint}{ex:stirling1}
$P(K_n=k)=|s(n,k)|/n!$. Apply Theorem~\ref{thm:modpoissonlocal} with $\lambda_n=\log n$ and $\psi=\frac1\Gamma$.
\end{hint}

\section{Solutions}

\begin{solution}{ex:shift}
\idea The Poisson law is its own reference; an independent addition multiplies generating functions.
\writeup For $X_n$ Poisson$(\lambda_n)$, $e^{-\lambda_n(z-1)}\mathbb Ez^{X_n}=1$ for all $z$. If $Y$ is independent, bounded and integer-valued, $\mathbb Ez^{X_n+Y}=\mathbb Ez^{X_n}\mathbb Ez^Y$, so $\psi_n(z)=\mathbb Ez^Y$, a polynomial, continuous with value $1$ at $1$.
Both $X_n$ and $X_n+Y$ satisfy the central limit theorem with the same normalization. But Theorem~\ref{thm:modpoissonlocal} says $P(X_n+Y=k)\approx\pi_{\lambda_n}(k)\,\mathbb E[x^Y]$ at $k=x\lambda_n$: for $x=2$ and $Y=1$ a.s., the probability is doubled, as it should be,
since $\pi_\lambda(k-1)=\frac k\lambda\pi_\lambda(k)$.
\end{solution}

\begin{solution}{ex:lambdas}
\idea The normalizing constant $\zeta(s)$ blows up as $s\downarrow1$, so every bounded range of integers loses its mass; and $\log\zeta(s)$ is essentially the sum over primes.
\writeup $P(N_s\le y)=\frac{\sum_{n\le y}n^{-s}}{\zeta(s)}\le(s-1)\sum_{n\le y}\frac1n\to0$. By Proposition~\ref{prop:dirichletint}, $\mathbb E\,\omega(N_s)=\sum_pP(p\mid N_s)=\lambda_s$. From Euler's product,
$\log\zeta(s)=\sum_p-\log(1-p^{-s})=\lambda_s+\sum_p\sum_{m\ge2}\frac{p^{-ms}}m$, and the double sum is at most $\sum_p\frac{p^{-2}}{1-p^{-1}}\le\sum_p\frac2{p^2}<\infty$. Finally $\frac1{s-1}\le\zeta(s)\le1+\frac1{s-1}$ (compare the sum with $\int_1^\infty t^{-s}\dd t$),
so $\log\zeta(s)=\log\frac1{s-1}+O(s-1)$, and $\lambda_s=\log\frac1{s-1}+O(1)$.
\end{solution}

\begin{solution}{ex:residues}
\idea Both sums have the same variances, so the same Gaussian part; their fourth cumulants differ, and the residue records the difference.
\writeup $\mathbb Ee^{itY_n}e^{H_nt^2/2}=\prod_{k\le n}\frac{\sin(\sqrt3t/\sqrt k)}{\sqrt3t/\sqrt k}e^{t^2/2k}$. With $v=\sqrt3t/\sqrt k$, $\log\frac{\sin v}v+\frac{v^2}6=-\frac{v^4}{180}+O(v^6)=-\frac{t^4}{20k^2}+O(\frac{t^6}{k^3})$, so as in
Example~\ref{ex:harmonicwalk} the product converges locally uniformly to $\psi_U(t)=\prod_k\frac{\sin(\sqrt3t/\sqrt k)}{\sqrt3t/\sqrt k}e^{t^2/2k}$. Then
\[
\begin{gathered}
\log\psi_U(t)=-\frac{t^4}{20}\sum_k\frac1{k^2}+O(t^6)=-\frac{\pi^2t^4}{120}+O(t^6),\\\log\psi(t)=-\frac{t^4}{12}\sum_k\frac1{k^2}+O(t^6)=-\frac{\pi^2t^4}{72}+O(t^6)
\end{gathered}
\]
for the signs. The coefficients differ, so $\psi_U\neq\psi$, although both sequences satisfy the central limit theorem with variance $H_n$. (The fourth cumulant of a sign is $-2$, that of $U$ is $-\frac65$; the residue sees $\sum_k\frac{\kappa_4}{24k^2}t^4$.)
\end{solution}

\begin{solution}{ex:bigomega}
\idea With multiplicity, each prime contributes a geometric variable, whose generating function has a pole; the smallest prime gives the nearest pole.
\writeup $P(p^m\mid N_s)=p^{-ms}$, and as in Proposition~\ref{prop:dirichletint} the exponents $v_p(N_s)$ are independent (the product rule holds for the events $\{p_i^{m_i}\mid N_s\}$). So $v_p$ is geometric: $P(v_p=m)=p^{-ms}(1-p^{-s})$, with
$\mathbb Ez^{v_p}=\frac{1-p^{-s}}{1-zp^{-s}}$ for $|z|<p^s$, and $\mathbb Ez^{\Omega(N_s)}=\prod_p\frac{1-p^{-s}}{1-zp^{-s}}$ for $|z|<2^s$. Then $e^{-\lambda_s(z-1)}\mathbb Ez^{\Omega}=\prod_p\frac{1-p^{-s}}{1-zp^{-s}}e^{-(z-1)p^{-s}}$, and
$\log\frac{1-u}{1-zu}-(z-1)u=O(|u|^2)$ uniformly for $|z|\le b<2$ and $u=p^{-s}\le\frac12$ with $|zu|\le\frac b2<1$; as in Theorem~\ref{thm:omegamod}, the product converges uniformly on $|z|\le b$ to the stated limit. The annulus cannot reach $|z|=2$: the factor
$\frac{1-2^{-s}}{1-z2^{-s}}$ has a pole at $z=2^s\to2$, and indeed $\mathbb E[2^{\Omega(N_s)}]=\prod_p\frac{1-p^{-s}}{1-2p^{-s}}$, finite for every $s>1$, tends to infinity as $s\downarrow1$ (the factor $p=2$ alone is $\sum_m2^m2^{-ms}(1-2^{-s})=\frac{1-2^{-s}}{1-2^{1-s}}$, which blows up). Integers with $k>2\lambda$ prime
factors are dominated by powers of $2$ times typical integers, a regime where the Poisson approximation, corrected or not, is the wrong model.
\end{solution}

\begin{solution}{ex:stirling1}
\idea The Stirling numbers are $n!$ times the probabilities of the number of cycles; the local theorem with residue $1/\Gamma$ gives them.
\writeup By Theorem~\ref{thm:cyclesmod} with parameters $\lambda=\log n$ and $\psi=\frac1\Gamma$, Theorem~\ref{thm:modpoissonlocal} gives, uniformly for $x=\frac k{\log n}\in[a,b]$,
\[
\frac{|s(n,k)|}{n!}=P(K_n=k)\sim\frac1n\cdot\frac{(\log n)^k}{k!}\cdot\frac1{\Gamma(x)}
\]
(since $e^{-\log n}=\frac1n$, and $\frac1{\Gamma(x)}$ is bounded away from $0$ on $[a,b]$, so the $o(1)$ becomes a relative error). Hence
\[
|s(n,k)|\sim\frac{(n-1)!\,(\log n)^k}{k!\,\Gamma(x)}=\frac{(n-1)!\,(\log n)^{k-1}}{(k-1)!}\cdot\frac{\log n}{k\,\Gamma(x)}=\frac{(n-1)!\,(\log n)^{k-1}}{(k-1)!\,\Gamma(1+x)},
\]
using $x\Gamma(x)=\Gamma(1+x)$. Leo Moser and Max Wyman studied these asymptotics by the saddle point method in 1958; this uniform form, with a full expansion, is due to Hsien-Kuei Hwang (1995).
\end{solution}

\section*{Sources and further reading}

\begin{itemize}
\item \emph{Mod-$\varphi$ convergence.} J.~P.~Keating and N.~C.~Snaith, ``Random matrix theory and $\zeta(1/2+it)$'', \emph{Comm.\ Math.\ Phys.}\ 214 (2000), 57--89. J.~Jacod, E.~Kowalski and A.~Nikeghbali, ``Mod-Gaussian convergence: new limit theorems
in probability and number theory'', \emph{Forum Math.}\ 23 (2011), 835--873. E.~Kowalski and A.~Nikeghbali, ``Mod-Poisson convergence in probability and number theory'', \emph{Int.\ Math.\ Res.\ Not.}\ 2010, 3549--3587, where the factorization of the residue
of $\omega(n)$ into $\frac1\Gamma$ times an Euler product is explained. V.~F\'eray, P.-L.~M\'eliot and A.~Nikeghbali, \emph{Mod-$\phi$ Convergence: Normality Zones and Precise Deviations}, SpringerBriefs, Springer, 2016. H.-K.~Hwang, ``On convergence rates in the central limit
theorems for combinatorial structures'', \emph{European J.\ Combin.}\ 19 (1998), 329--343, the ``quasi-powers'' theorem, an ancestor of Theorem~\ref{thm:modpoissonlocal}.
\item \emph{Cycles and prime factors.} V.~Gontcharoff, ``Du domaine de l'analyse combinatoire'', \emph{Izv.\ Akad.\ Nauk SSSR Ser.\ Mat.}\ 8 (1944), 3--48. L.~Moser and M.~Wyman, ``Asymptotic development of the Stirling numbers of the first kind'', \emph{J.\ London Math.\ Soc.}\ 33 (1958), 133--146. H.-K.~Hwang, ``Asymptotic expansions for the Stirling numbers of the first kind'', \emph{J.\ Combin.\ Theory Ser.\ A} 71 (1995), 343--351.
A.~Selberg, ``Note on a paper by L.~G.~Sathe'', \emph{J.\ Indian Math.\ Soc.}\ 18 (1954), 83--87. H.~Delange, ``Sur des formules de Atle Selberg'', \emph{Acta Arith.}\ 19 (1971), 105--146. A.~R\'enyi and P.~Tur\'an, ``On a theorem of Erd\H os--Kac'', \emph{Acta Arith.}\ 4 (1958), 71--84.
\item G.~Tenenbaum, \emph{Introduction to Analytic and Probabilistic Number Theory}, 3rd ed., AMS, 2015, Chapter~II.5: the Selberg--Delange method in full.
\item R.~Arratia, A.~D.~Barbour and S.~Tavar\'e, \emph{Logarithmic Combinatorial Structures: a Probabilistic Approach}, EMS, 2003. Permutations, polynomials over finite fields and integers, side by side.
\end{itemize}
